% Layout C — loaded by styles/versatus-covers.sty.
% Select it with:  \renewcommand{\VSCoverOption}{LayoutC}
% Defines \VSCoverLayoutC. This is NOT a package: the package owns the preamble.



% --- Comando Principal da Capa ---
\newcommand{\VSCoverLayoutC}{%
\begin{titlepage}
\thispagestyle{empty}
\noindent
\begin{tikzpicture}[remember picture, overlay, shift={(current page.south west)}, x=1cm, y=1cm]
  
  % 1. Fundo Base
  \fill[VSPaperWhite] (-1, -1) rectangle (22, 31);
  
  % 2. BLOCO SUPERIOR DE LOSANGOS (Iteração Matemática)
  \foreach \row in {0, 1, 2, 3, 4} {
    \foreach \col in {0, 1, 2, 3, 4, 5} {
      \pgfmathsetmacro{\cx}{-1.25 + \col * 5 + mod(\row,2)*2.5}
      \pgfmathsetmacro{\cy}{32 - \row * 2.5}
      
      % Alterna entre as 3 cores baseadas na posição do grid
      \pgfmathparse{int(mod(\col+\row, 3))}
      \ifnum\pgfmathresult=0 \def\dcol{VSTeal} \fi
      \ifnum\pgfmathresult=1 \def\dcol{VSBlack} \fi
      \ifnum\pgfmathresult=2 \def\dcol{VSRed} \fi
      
      \fill[\dcol] (\cx, \cy+2.5) -- (\cx+2.5, \cy) -- (\cx, \cy-2.5) -- (\cx-2.5, \cy) -- cycle;
    }
  }

  % 3. BLOCO INFERIOR DE LOSANGOS
  \foreach \row in {0, 1, 2, 3, 4} {
    \foreach \col in {0, 1, 2, 3, 4, 5} {
      \pgfmathsetmacro{\cx}{-1.25 + \col * 5 + mod(\row,2)*2.5}
      \pgfmathsetmacro{\cy}{7 - \row * 2.5}
      
      % Inverte a ordem das cores para criar movimento
      \pgfmathparse{int(mod(\col+\row, 3))}
      \ifnum\pgfmathresult=0 \def\dcol{VSRed} \fi
      \ifnum\pgfmathresult=1 \def\dcol{VSTeal} \fi
      \ifnum\pgfmathresult=2 \def\dcol{VSBlack} \fi
      
      \fill[\dcol] (\cx, \cy+2.5) -- (\cx+2.5, \cy) -- (\cx, \cy-2.5) -- (\cx-2.5, \cy) -- cycle;
    }
  }

  % 4. FAIXA CENTRAL (Área de Respiro para o Texto)
  \fill[VSPaperWhite] (-1, 10) rectangle (22, 19.5);
  
  % 5. Logomarca na faixa central
  \node[anchor=north west, inner sep=0] at (11.5, 29) {
    \VSLogoEscolhida[1.5] % O 1.5 é a escala. Pode ajustar para 1.2 ou 1.8 se quiser maior/menor
};
  
  % 6. Título Principal (Estilo Vision)
  \node[anchor=west, text=VSBlack, inner sep=0, text width=12cm] at (1.5, 14.5) {
    \linespread{0.8}\selectfont
    \fontsize{65}{65}\selectfont\bfseries{\BookTitle}\par
  };
  
  % 7. Dados adicionais estruturados à direita
  \node[anchor=west, text=VSBlack, inner sep=0, text width=6cm] at (14, 14.5) {
    {\fontsize{10}{12}\selectfont\textbf{\BookAuthor}\par}
    \vspace{0.1cm}
    {\fontsize{10}{12}\selectfont\BookDate\par}
    \vspace{0.8cm}
    {\fontsize{10}{12}\selectfont\textbf{\BookSubtitle}\par}
    \vspace{0.1cm}
    {\fontsize{10}{12}\selectfont\color{VSMutedText}\BookDescription\par}
  };

\end{tikzpicture}
\end{titlepage}%
\clearpage
}
